\documentclass[11pt,letterpaper]{article}
\usepackage[T1]{fontenc}
\usepackage[utf8]{inputenc}
\usepackage{amsmath,amsthm,amssymb,mathtools}
\usepackage{newpxtext,newpxmath}
\usepackage[margin=0.95in,headheight=15pt]{geometry}
\usepackage{microtype}
\usepackage{booktabs,tabularx,longtable,array}
\usepackage{xcolor,enumitem,fancyhdr,needspace}
\usepackage{titlesec}
\usepackage[unicode=true,colorlinks=true,linkcolor=blue!45!black,citecolor=blue!45!black,urlcolor=blue!45!black,breaklinks=true]{hyperref}
\usepackage{bookmark}
\definecolor{ink}{HTML}{173C4B}
\titleformat{\section}{\Large\bfseries\color{ink}}{\thesection}{0.7em}{}
\titleformat{\subsection}{\large\bfseries\color{ink}}{\thesubsection}{0.7em}{}
\titlespacing*{\section}{0pt}{2.2ex plus 1ex minus .2ex}{1.0ex plus .2ex}
\titlespacing*{\subsection}{0pt}{1.8ex plus 1ex minus .2ex}{.7ex plus .2ex}
\setlength{\parskip}{0.35em}
\setlength{\parindent}{1.1em}
\setlength{\emergencystretch}{2em}
\setlist{itemsep=0.25em,topsep=0.4em,leftmargin=1.5em}
\pagestyle{fancy}
\fancyhf{}
\fancyhead[L]{\small Bounded-support Hahn arithmetic}
\fancyhead[R]{\small Research draft}
\fancyfoot[C]{\thepage}
\renewcommand{\headrulewidth}{0.3pt}
\newtheorem{theorem}{Theorem}[section]
\newtheorem{lemma}[theorem]{Lemma}
\newtheorem{proposition}[theorem]{Proposition}
\newtheorem{corollary}[theorem]{Corollary}
\theoremstyle{definition}
\newtheorem{definition}[theorem]{Definition}
\newtheorem{example}[theorem]{Example}
\newtheorem{construction}[theorem]{Construction}
\theoremstyle{remark}
\newtheorem{remark}[theorem]{Remark}
\newtheorem{warning}[theorem]{Warning}
\newcommand{\Q}{\mathbb Q}
\newcommand{\R}{\mathbb R}
\newcommand{\C}{\mathbb C}
\newcommand{\N}{\mathbb N}
\newcommand{\Z}{\mathbb Z}
\newcommand{\No}{\mathbf{No}}
\newcommand{\On}{\mathbf{On}}
\newcommand{\supp}{\operatorname{supp}}
\newcommand{\Frac}{\operatorname{Frac}}
\newcommand{\trdeg}{\operatorname{trdeg}}
\newcommand{\cof}{\operatorname{cf}}
\newcommand{\lc}{\operatorname{lc}}
\newcommand{\Aut}{\operatorname{Aut}}
\newcommand{\Span}{\operatorname{span}}
\newcommand{\HH}[2]{#1((t^{#2}))}
\newcommand{\BB}[2]{\mathcal B_{#2}(#1)}
\newcommand{\FF}[2]{\mathcal F_{#2}(#1)}
\newcommand{\cc}{\mathfrak c}
\newcommand{\abs}[1]{\lvert #1\rvert}
\newcommand{\set}[1]{\{#1\}}
\newcommand{\repo}{\texttt{VladimirReshetnikov/Surreal}}
\newcommand{\auditcommit}{\texttt{048b72cf7cbfc8ab246e4f73788c10460cb3f6e0}}
\hypersetup{pdftitle={Bounded-Support Hahn Arithmetic: Cofinal Descent and Maximal Transcendence},pdfauthor={Research draft prepared with ChatGPT},pdfsubject={Surreal and surcomplex numbers; Hahn fields; algebraic independence; cofinality}}
\begin{document}
\begin{titlepage}
\thispagestyle{empty}
\vspace*{0.4in}
{\small\sffamily\bfseries\color{ink} SURREAL AND SURCOMPLEX RESEARCH}\par
\vspace{0.25in}
{\Huge\bfseries\color{ink} Bounded-Support\par Hahn Arithmetic}\par
\vspace{0.2in}
{\LARGE Cofinal descent, optimal support thresholds,\par and maximal transcendence}\par
\vspace{0.3in}
{\large Research draft prepared with ChatGPT\par 22 September 2026}\par
\vspace{0.25in}
\noindent\rule{\textwidth}{0.6pt}
\begin{abstract}
Let $G$ be a nonzero set-sized ordered abelian group, let $K$ have characteristic
zero, and let $\mathcal B_G(K)$ be the ring of Hahn series whose support is
bounded above in $G$. Put $\mathcal F_G(K)=\Frac\mathcal B_G(K)$ and
$\kappa=\cof(G)$. We construct $2^\kappa$ elements of $K((t^G))$, all with
positive integer coefficients and one common support of order type $\kappa$,
that are algebraically independent over $\mathcal F_G(K)$. No smaller support
cardinality can support even one transcendental element over this base.
Every valuation ball contains such an independent family. In particular, the
transcendence degree is the continuum for real-exponent Hahn fields over
$\R$ or $\C$.

A complementary descent theorem proves that, whenever $H\subseteq G$ is
cofinal and $K\subseteq L$, the fields $K((t^H))$ and $\mathcal F_G(L)$ are
linearly disjoint over $\mathcal F_H(K)$. Polynomial relation ideals and finite
algebraic degrees therefore survive simultaneous enlargement of coefficients
and cofinal exponent groups. We also classify when $\mathcal B_G(K)$ is a
field, reduce its unit theory to the highest Archimedean quotient, and prove
a conditional GCD-descent theorem relevant to a published pre-Schreier
question. Concrete Conway-normal-form specializations give independent
families of actual surreal numbers, including maximal families of size
$2^\kappa$ at every infinite regular cardinal.
\end{abstract}
\vfill
\noindent\small\textbf{Status.} Complete arguments for the stated new results are given relative
to explicitly identified standard Hahn-field and Conway-normal-form foundations.
The proposed novelty is not a certified priority claim. This draft is not refereed
or Lean-verified. The GCD application is explicitly conditional and is not used
in any transcendence theorem.
\end{titlepage}
\begingroup
\small
\setlength{\parskip}{0pt}
\renewcommand{\baselinestretch}{0.95}\selectfont
\tableofcontents
\endgroup
\clearpage

\section{The problem, the results, and their provenance}\label{sec:intro}

A Hahn series can contain infinitely many terms while its exponents remain in
one bounded interval. Such a series is not a polynomial, and adjoining all
such series produces a much larger base than a rational function field.
The question pursued here is what remains genuinely transcendental after
\emph{every} bounded-support series has been admitted as a coefficient, and
how that answer changes when the ambient exponent group changes.

Write
\begin{equation}\label{eq:base}
 \BB KG=\{f\in\HH KG:\supp(f)\text{ is bounded above in }G\},
 \qquad \FF KG=\Frac\BB KG.
\end{equation}
The phrase ``bounded support'' always refers to the order on the specified
exponent group. It does not mean finite support, bounded coefficients, or a
bound on the cardinality of the support. In particular it is not the
cardinal restriction sometimes called a $\lambda$-bounded Hahn field.

\subsection{The principal theorems}

Here is the theorem package in a form that can be read before its proofs.
All fields and groups in these statements are sets. Cardinal assertions are
made in set theory with choice.

\begin{theorem}[Optimal-support independence; proved in Section~\ref{sec:independence}]
\label{thm:mainintro}
Let $K$ have characteristic zero and $G\ne0$ be an ordered abelian group.
Set $\kappa=\cof(G)$. There are a strictly increasing cofinal sequence
$(a_\alpha)_{\alpha<\kappa}$ of positive elements of $G$ and positive integers
$c_A(\alpha)$, for $A\subseteq\kappa$, such that
\[
 Y_A=\sum_{\alpha<\kappa}c_A(\alpha)t^{a_\alpha}
 \quad(A\subseteq\kappa)
\]
are algebraically independent over $\FF KG$. Their common support has order
type $\kappa$. The minimum possible cardinality, and the minimum possible
order type, of the support of a series transcendental over $\FF KG$ are both
$\kappa$. Every nonempty valuation ball contains a family of $2^\kappa$
elements algebraically independent over $\FF KG$.
\end{theorem}

The strength is not a cardinality comparison between the base and the larger
field: they can have the same cardinality. It is the construction of a free
polynomial algebra over the entire bounded-support fraction field, with
integer-coefficient witnesses and an optimal support requirement.

\begin{theorem}[Cofinal coefficient--exponent descent; proved in Section~\ref{sec:descent}]
\label{thm:descentintro}
Let $K\subseteq L$ be fields and let $0\ne H\subseteq G$ be ordered abelian
groups. If $H$ is cofinal in $G$, multiplication induces an injection
\[
 \FF LG\otimes_{\FF KH}\HH KH\ \longrightarrow\ \HH LG.
\]
Consequently
\[
 \FF LG\cap\HH KH=\FF KH.
\]
Every finite tuple from $\HH KH$ has the same polynomial relation ideal
after extension of scalars from $\FF KH$ to $\FF LG$. Finite algebraic
degrees are preserved. In characteristic zero the cofinality hypothesis is
also necessary for the displayed multiplication map to be injective.
\end{theorem}

For example, with $H=\Z$, $G=\R$, $K=\Q$, and $L=\C$, every algebraic
relation among ordinary rational-coefficient Laurent series over
$\FF{\C}{\R}$ already comes from relations over $\Q(t)$. This is a
simultaneous coefficient and exponent extension statement, not merely an
intersection statement.

The structural companion is an exact distinction between the existence and
absence of a highest scale. The ring $\BB KG$ is a field precisely when $G$
has no order unit. In that case it is the union of the Hahn fields over
proper convex subgroups. For divisible $G$ and real closed or algebraically
closed $K$, that union is itself real closed or algebraically closed.
When an order unit exists, the units are controlled by the highest
Archimedean quotient; see Section~\ref{sec:rank}.

\subsection{What is already known, and what is proposed here}

L'Innocente and Mantova study bounded-support localization and its arithmetic
in \cite{LM24}. Their Proposition~2.4.5 already gives a cofinality obstruction
to representing all Hahn series by fractions of one-sided series. That
nonrationality obstruction is \emph{not} claimed here as new. The proposed
strengthening is the $2^{\cof(G)}$-sized algebraically independent family with
one optimal common support, together with the coefficient--exponent descent
and its exact converse. Their Proposition~3.5.1 supplies the upper-support
multiplicativity used only in the unit and GCD discussion.

The repository audit used the documentation catalogue and the closest
transcendence report at revision
\begin{center}\small\auditcommit.\end{center}
The catalogue contains 36 reports. Its tail-span report \cite{RepoTail}
constructs differential independence using independent square classes in
coefficient extensions. Here the coefficients of the witnesses can all be
positive integers, the base is a bounded-support fraction field, and the
mechanism is support separation. These are different base fields and
different assertions; neither report should be described as automatically
subsuming the other. The catalogue also contains entire-function and
cofinality results, but a ring of entire functions is not the ring in
\eqref{eq:base}.

The proposed priority claims are deliberately limited. Targeted searches did
not identify the combined descent theorem or the sharp independent-family
theorem in these formulations. This is not an exhaustive MathSciNet,
Zentralblatt, thesis, or citation-network audit. The elementary coding
lemma, standard linear-disjointness algebra, Hahn field construction,
real/algebraic closedness criteria, and Conway normal form are background,
not new discoveries. A later repository revision, \texttt{6e34651fe640},
was observed during the audit; its recorded change concerned Hahn--Hilbert
spectral theory. No claim of a full source-by-source audit of that later
revision is made.

\subsection{Dependencies and the boundary of the claim}

The main independence proof uses only Hahn multiplication, finite polynomial
algebra, cardinal cofinality, and a fully proved finite-pattern coding
lemma. It does not use a GCD theorem, the upper-support multiplicativity
result, any surreal derivation, or analytic convergence. The descent proof
is coefficientwise and works in arbitrary characteristic. Standard
Hahn-field closedness and Conway normal form are invoked only where named.
Appendix~\ref{sec:gcd} proves an implication toward a published GCD/pre-Schreier
problem; it does not supply the missing real-exponent GCD input.

\section{Hahn series and bounded-support algebra}\label{sec:setup}

\subsection{Conventions}

An ordered abelian group is written additively. A Hahn series over $K$ is a
formal expression
\[
 f=\sum_{g\in G} f_g t^g
\]
whose support $\{g:f_g\ne0\}$ is well ordered. Addition is coefficientwise
and multiplication is the Hahn convolution. At each exponent the relevant
sum is finite. Finite unions and finite sumsets of well-ordered supports
are well ordered. These standard facts give the field $\HH KG$; see
\cite[Section~2.1]{LM24}. We use the conventional valuation
\[
 v(f)=\min\supp(f)\quad(f\ne0),\qquad v(0)=+\infty.
\]
Thus $v(fg)=v(f)+v(g)$ and $v(f+g)\ge\min(v(f),v(g))$.
The convention at zero is stated explicitly because some treatments of
one-sided series use a different special-symbol convention.

A family of Hahn series is \emph{strongly summable} when its combined support
is well ordered and only finitely many summands contribute at any exponent.
A sum along a strictly increasing sequence of exponents is strongly
summable, independently of whether its finite partial sums converge in a
chosen topology. All transfinite sums below are Hahn sums in this sense.

For $f\in\HH KG$ and $\delta\in G$, write
\[
 f_{\le\delta}=\sum_{g\le\delta} f_g t^g.
\]
A nonzero Hahn support has a least element. Therefore ``bounded above'' in
\eqref{eq:base} is equivalent to containment in some order interval
$[\ell,u]\subseteq G$.

\begin{lemma}[Elementary support bounds]\label{lem:bounds}
The set $\BB KG$ is a subring of $\HH KG$, contains $K$ and every monomial,
and is closed under taking arbitrary subseries. If
$\supp(f)\subseteq[\ell,u]$ and
$\supp(g)\subseteq[\ell',u']$, then
\[
 \supp(fg)\subseteq[\ell+\ell',u+u'].
\]
Every finite family of bounded-support series has common lower and upper
support bounds.
\end{lemma}
\begin{proof}
Every exponent of a product is a sum of one exponent from each factor.
Cancellation only removes terms and cannot create an exponent outside the
indicated interval. A finite union is bounded by the minimum of its lower
bounds and the maximum of its upper bounds. Negation leaves the support
unchanged. Every subset of a well-ordered set is well ordered, proving the
subseries assertion. These observations also cover zero by using arbitrary
bounds for its empty support.
\end{proof}

\begin{proposition}[Localization and density]\label{prop:localdensity}
For any nonzero $G$,
\[
 \BB KG=t^G K((t^{G^{\le0}})).
\]
Both $\BB KG$ and $\FF KG$ are dense in $\HH KG$ for the valuation topology.
\end{proposition}
\begin{proof}
If $\supp(f)\le u$, then $t^{-u}f$ has nonpositive support. Conversely a
monomial translate of a nonpositive support is bounded above. For density,
$f_{\le\delta}\in\BB KG$ and either $f=f_{\le\delta}$ or
$v(f-f_{\le\delta})>\delta$. Hence every basic valuation ball contains a
bounded-support element. The assertion for the fraction field follows.
\end{proof}

There is no assertion here that bounded-support series are finite sums.
For example
\[
 \sum_{n\ge2} t^{1-1/n}\in\BB{\R}{\Q}
\]
has infinitely many exponents in $[1/2,1)$. Nor is density an algebraicity
statement. The main theorem will show that an enormous transcendental
extension can remain after adjoining this dense subring's entire fraction
field.

\subsection{Cofinality}

\begin{definition}
For a nonzero ordered abelian group $G$, its cofinality $\cof(G)$ is the least
cardinality of a subset unbounded above in $G$. An ordered subgroup $H$ is
\emph{cofinal} if for every $g\in G$ some $h\in H$ satisfies $h>g$.
An \emph{order unit} is an element $u>0$ such that for every $g\in G$ some
positive integer $n$ satisfies $\abs g\le nu$.
\end{definition}

\begin{lemma}\label{lem:cofinality}
The cardinal $\kappa=\cof(G)$ is infinite and regular. There is a strictly
increasing positive cofinal sequence $(g_\alpha)_{\alpha<\kappa}$.
Every subset of $G$ of cardinality less than $\kappa$ is bounded above.
An ordered subgroup is either cofinal or bounded above.
\end{lemma}
\begin{proof}
No finite subset is cofinal, since it has a maximum and $G$ has positive
elements. The last two assertions follow from the definitions. Start with
a cofinal family of size $\kappa$ and recursively choose a value exceeding
its next member and all previously chosen values. At stage $\alpha<\kappa$
the latter set has size less than $\kappa$, so a bound exists. This gives
the increasing sequence. If $\kappa$ were singular, a cofinal subsequence
of its indices of cardinality $\cof(\kappa)<\kappa$ would give a smaller
cofinal subset of $G$, a contradiction.
\end{proof}

For a nonzero subgroup of $\R$, every positive element is an order unit and
the cofinality is $\aleph_0$. The converse implication from countable
cofinality to the existence of an order unit is false: a countable ascending
union of strictly increasing Archimedean scales has no highest scale.
Section~\ref{sec:surreal} gives concrete examples inside $\No$.

\section{Cofinal coefficient--exponent descent}\label{sec:descent}

The following result is useful independently of the lacunary constructions.
It explains why a bounded-support coefficient cannot conceal a new relation
among series supported on a cofinal subgroup.

\subsection{Coset projections and coefficient functionals}

Let $K\subseteq L$ be fields, $H\subseteq G$ an ordered subgroup, and
$r\in G$. Define the shifted coset projection
\begin{equation}\label{eq:projection}
 \pi_r(f)=\sum_{h\in H}f_{r+h}t^h\in\HH LH.
\end{equation}
It is defined for every $f\in\HH LG$: a subset of its support remains well
ordered, and translation preserves that property. This is not asserted to
be a ring homomorphism.

\begin{lemma}[Projection identity]\label{lem:projection}
For $z\in\HH LH$ and $b\in\HH LG$,
\[
 \pi_r(bz)=\pi_r(b)z.
\]
If $H$ is cofinal and $b\in\BB LG$, then $\pi_r(b)\in\BB LH$.
\end{lemma}
\begin{proof}
A term of $bz$ at an exponent in $r+H$ can only use a term of $b$ in that
same coset, because all exponents of $z$ lie in $H$. At every exponent the
convolution is finite, so selecting the coset commutes with it. If the
support of $b$ is bounded above by $u\in G$, then the shifted support is
bounded above in $G$ by $u-r$. Cofinality supplies an $H$-valued bound.
\end{proof}

Let $\lambda:L\to K$ be any $K$-linear functional. Apply it to each
coefficient to obtain
\[
 \lambda_*:\HH LH\longrightarrow\HH KH.
\]
This map can kill support terms but introduces none. In particular it
preserves boundedness of support. For $z\in\HH KH$ it satisfies
\begin{equation}\label{eq:lambda}
 \lambda_*(bz)=\lambda_*(b)z.
\end{equation}
The reason is again the finiteness of each convolution coefficient, now
combined with $K$-linearity. No multiplicativity of $\lambda$ is needed.
For any nonzero $a\in L$, a $K$-linear functional with $\lambda(a)\ne0$
exists: extend $a$ to a basis and prescribe its value to be $1$.

\subsection{The descent theorem}

\begin{theorem}[Simultaneous descent]\label{thm:descent}
Let $K\subseteq L$ be fields and let $0\ne H\subseteq G$ be cofinal.
Every finite family of elements of $\HH KH$ linearly independent over
$\FF KH$ remains linearly independent over $\FF LG$. Equivalently, the
natural multiplication map
\begin{equation}\label{eq:tensor}
 \FF LG\otimes_{\FF KH}\HH KH\longrightarrow\HH LG
\end{equation}
is injective.
\end{theorem}
\begin{proof}
Let $z_1,\ldots,z_m\in\HH KH$ be linearly independent over $\FF KH$.
Suppose there is a nonzero linear relation over $\FF LG$. After multiplying
by a common nonzero denominator, it becomes
\begin{equation}\label{eq:linearrel}
 \sum_{j=1}^m b_jz_j=0,\qquad b_j\in\BB LG,
\end{equation}
with some $b_{j_0}\ne0$. Choose an exponent $r\in\supp(b_{j_0})$ and set
$a=(b_{j_0})_r\ne0$. Applying $\pi_r$ to \eqref{eq:linearrel} gives
\[
 \sum_j\pi_r(b_j)z_j=0,
\]
where all projected coefficients belong to $\BB LH$. Choose a $K$-linear
functional $\lambda:L\to K$ with $\lambda(a)=1$. Applying $\lambda_*$ and
using \eqref{eq:lambda} yields
\[
 \sum_j\lambda_*(\pi_r(b_j))z_j=0.
\]
Its coefficients belong to $\BB KH\subseteq\FF KH$. The coefficient of
$t^0$ in $\lambda_*(\pi_r(b_{j_0}))$ is $1$, so the relation is nonzero.
This contradicts the assumed independence.

To verify the tensor formulation directly, write a tensor using a finite
$\FF KH$-linearly independent family of second factors. If its image is
zero, the first part of the proof forces every first-factor coefficient to
vanish. The converse equivalence follows by the same representation.
This agrees with the standard definition of linear disjointness
\cite{StacksLD}.
\end{proof}

\begin{corollary}[Intersection, relations, and algebraic degree]\label{cor:relations}
Under the hypotheses of Theorem~\ref{thm:descent}:
\begin{enumerate}
\item $\FF LG\cap\HH KH=\FF KH$ inside $\HH LG$.
\item For a finite tuple $z$ in $\HH KH$, let $I_K(z)$ and $I_L(z)$ be its
relation ideals over $\FF KH$ and $\FF LG$, respectively. Then
\[
 I_L(z)=I_K(z)\,\FF LG[X_1,\ldots,X_m].
\]
\item Algebraic independence of any family in $\HH KH$ is preserved in both
directions by this base enlargement.
\item If $z\in\HH KH$ is algebraic of degree $d$ over $\FF KH$, it has degree
$d$ over $\FF LG$, with the same monic minimal polynomial.
\end{enumerate}
\end{corollary}
\begin{proof}
For the intersection statement, $z\notin\FF KH$ makes $1,z$ linearly
independent over $\FF KH$, hence over $\FF LG$; it cannot then belong to
$\FF LG$. The reverse inclusion is immediate.

Let $E=\FF KH$ and $F=\FF LG$. The quotient $E[X]/I_K(z)$ embeds in
$\HH KH$. Tensoring that inclusion with the $E$-vector space $F$ preserves
injectivity. By Theorem~\ref{thm:descent}, the resulting algebra still embeds
in $\HH LG$. Its presentation is $F[X]/I_K(z)F[X]$, so its evaluation
kernel is exactly the asserted ideal. This proves the second assertion and
then the third. Finally $1,z,\ldots,z^{d-1}$ remain linearly independent,
while the old degree-$d$ equation still holds, proving the last assertion.
\end{proof}

For a finite Galois subextension of $\HH KH/\FF KH$, the same argument
preserves the degree and Galois group after compositum with $\FF LG$.
Indeed, the tensor product is a field of the original degree; its distinct
base-field automorphisms extend by the identity on $\FF LG$, and their
number already equals that degree. This observation does not assert that
every algebraic germ extension is Galois.

\subsection{The ordinary Laurent-series case}

\begin{corollary}[One cofinal scale]\label{cor:onescale}
If $u>0$ is an order unit of $G$ and $K\subseteq L$, then
$K((t^u))$ and $\FF LG$ are linearly disjoint over $K(t^u)$. In particular
\[
 K((t^u))\cap\FF LG=K(t^u).
\]
\end{corollary}
\begin{proof}
The subgroup $\Z u$ is cofinal. Its bounded well-ordered subsets are finite,
so $\BB K{\Z u}=K[t^u,t^{-u}]$ and its fraction field is $K(t^u)$.
Apply Theorem~\ref{thm:descent}.
\end{proof}

\begin{example}[An algebraic germ does not become a bounded fraction]
Take $s=t^u$ and characteristic zero. The binomial series
\[
 z=(1+s)^{1/2}=1+\tfrac12s-\tfrac18s^2+\tfrac1{16}s^3-\cdots
 \in\Q[[s]]
\]
has degree two over $\FF LG$ for every characteristic-zero extension
$\Q\subseteq L$ and every $G$ in which $u$ is an order unit.
To check its degree over $\Q(s)$, observe that $1+s$ is not a rational
square: in a coprime quotient of polynomials, the multiplicity of the
factor $s+1$ would be even on a square and odd on $1+s$. Corollary~\ref{cor:relations}
then preserves degree two. Thus it is algebraic but not itself in the
bounded-support fraction field. In contrast, $(1-s)^{-1}$ is in that
fraction field, although its support is unbounded.
\end{example}

\subsection{Why noncofinal enlargement behaves oppositely}

If $H$ is not cofinal, the entire group $H$ is bounded above in $G$.
Consequently
\begin{equation}\label{eq:collapse}
 \HH KH\subseteq\BB LG\subseteq\FF LG.
\end{equation}
Every old Hahn series becomes a bounded-support coefficient in the enlarged
workspace. Once Theorem~\ref{thm:independent} has established
$\FF KH\ne\HH KH$ in characteristic zero, this proves the converse in
Theorem~\ref{thm:descentintro}: choose $z$ in that difference. The pair
$1,z$ is independent over the smaller base but dependent over the larger
one. The exceptional subgroup $H=0$ is excluded because its Hahn field and
fraction field are both just $K$.

This dichotomy is about the actual embedding of exponent groups. Merely
saying that both groups are divisible or both have countable cofinality does
not imply that the embedding is cofinal.

\section{The support-separation mechanism}\label{sec:gaps}

This section supplies the estimate that upgrades the earlier nonrationality
obstruction to algebraic independence. It is entirely algebraic. Even when
a truncation has infinitely many terms, its support has the bound used in
the estimate.

\subsection{Separated exponent systems}

\begin{definition}\label{def:skeleton}
Let $\kappa$ be an infinite cardinal. A \emph{separated exponent system} in
$G$ consists of positive elements $a_\alpha,b_\alpha$ for $\alpha<\kappa$
such that:
\begin{enumerate}
\item $a_\alpha$ is strictly increasing and cofinal, and
$a_\beta\le b_\alpha$ whenever $\beta<\alpha$;
\item for every positive integer $d$ and every $w\in G$, eventually in
$\alpha$ one has
\begin{equation}\label{eq:gaps}
 a_\alpha>d b_\alpha+w,
 \qquad a_{\alpha+1}>d a_\alpha+w.
\end{equation}
\end{enumerate}
Every infinite cardinal is a limit ordinal, so $\alpha+1<\kappa$ for all
$\alpha<\kappa$.
\end{definition}

\begin{lemma}[Existence at the exact cofinality]\label{lem:skeleton}
Every nonzero ordered abelian group admits a separated exponent system of
length $\kappa=\cof(G)$.
\end{lemma}
\begin{proof}
Fix a strictly increasing positive cofinal sequence $g_\alpha$ from
Lemma~\ref{lem:cofinality}.

First suppose $\kappa=\aleph_0$, and index by $n\ge0$. Choose $b_n>0$
strictly exceeding $g_n$ and, when $n>0$, $a_{n-1}$. Set
$a_n=(n+2)b_n$. These elements are increasing and cofinal, and earlier
$a_j$ are at most $b_n$. Given $d,w$, eventually $b_n>w$ and $n+2\ge d+1$,
so $a_n-d b_n\ge b_n>w$. At the next index the same estimate gives
$a_{n+1}>d b_{n+1}+w\ge d a_n+w$. Both inequalities hold eventually.

Now suppose $\kappa$ is uncountable. At stage $\alpha$, the earlier
$a_\beta$, together with $g_\alpha$, form a set of size less than $\kappa$.
Choose a positive strict upper bound $b_\alpha$ for them. The countable set
$\{n b_\alpha:n\ge1\}$ also has size less than $\kappa$, so choose
$a_\alpha$ strictly above every one of its elements. This gives a strictly
increasing cofinal sequence. For fixed $d,w$, eventually $b_\alpha>w$;
then
$a_\alpha>(d+1)b_\alpha>d b_\alpha+w$.
The estimate at $\alpha+1$, and $b_{\alpha+1}>a_\alpha$, give the second
inequality of \eqref{eq:gaps}.
\end{proof}

\begin{example}[Factorial spacing]
If $u$ is an order unit, the concrete choices
\[
 a_n=n!u,\qquad b_n=(n-1)!u\qquad(n\ge2)
\]
form a separated exponent system. For fixed $d$,
$a_n-d b_n=(n-d)(n-1)!u$ eventually exceeds every $w\in G$.
The next-step inequality is checked in the same way.
\end{example}

\subsection{A polynomial cannot bridge two separated support bands}

Let
\[
 P(X)=\sum_{\nu}p_\nu X^\nu\in\BB KG[X_1,\ldots,X_m]
\]
be nonzero of total degree $d\ge1$. Its top homogeneous part is denoted
$P_d$. Choose $\ell,u\in G$ such that every $p_\nu$ is supported in
$[\ell,u]$. For coefficient vectors
$c(\alpha)=(c_1(\alpha),\ldots,c_m(\alpha))\in K^m$, put
\[
 Y_j=\sum_{\alpha<\kappa}c_j(\alpha)t^{a_\alpha}.
\]
Zero entries cause no problem: each support is a subset of the well-ordered
set of $a_\alpha$.

\begin{lemma}[Support-band detection]\label{lem:detection}
Suppose an index $\alpha$ satisfies
\begin{equation}\label{eq:twogaps}
 a_\alpha>d b_\alpha+(u-\ell),\qquad
 a_{\alpha+1}>d a_\alpha+(u-\ell),
\end{equation}
and $P_d(c(\alpha))\ne0$ in $\HH KG$. Then $P(Y_1,\ldots,Y_m)\ne0$.
\end{lemma}
\begin{proof}
Define the earlier part and the truncation through $\alpha$ by
\[
 A_j=\sum_{\beta<\alpha}c_j(\beta)t^{a_\beta},\qquad
 T_j=A_j+c_j(\alpha)t^{a_\alpha}.
\]
Every $A_j$ is supported in $[0,b_\alpha]$ and every $T_j$ in
$[0,a_\alpha]$. Expanding in an auxiliary \emph{polynomial} variable $Z$,
write
\begin{equation}\label{eq:auxpoly}
 P(A_1+c_1(\alpha)Z,\ldots,A_m+c_m(\alpha)Z)
   =\sum_{k=0}^d Q_k Z^k.
\end{equation}
The coefficient of $Z^d$ is exactly $P_d(c(\alpha))$, so $Q_d\ne0$.
Every other coefficient is a finite sum of a $p_\nu$, a scalar, and at most
$d$ factors chosen from the $A_j$. Lemma~\ref{lem:bounds} therefore gives
\[
 \supp(Q_k)\subseteq[\ell,u+d b_\alpha]\qquad(0\le k\le d).
\]
After substituting $Z=t^{a_\alpha}$, these supports lie in the bands
\begin{equation}\label{eq:bands}
 [\ell+k a_\alpha,\ u+d b_\alpha+k a_\alpha].
\end{equation}
The first inequality in \eqref{eq:twogaps} says that successive bands are
strictly disjoint. Since the last band has a nonzero contribution,
$P(T_1,\ldots,T_m)\ne0$.

A direct bound using the original polynomial and the supports of $T_j$
shows that
\begin{equation}\label{eq:truncbound}
 \supp(P(T))\subseteq[\ell,u+d a_\alpha].
\end{equation}
Meanwhile every tail $Y_j-T_j$ is either zero or has valuation at least
$a_{\alpha+1}$. Expand each difference $Y^\nu-T^\nu$ as a sum of products
containing at least one such tail. All other series factors have
nonnegative valuation. Thus
\begin{equation}\label{eq:tailbound}
 v(P(Y)-P(T))\ge\ell+a_{\alpha+1}.
\end{equation}
By the second inequality in \eqref{eq:twogaps}, this lower bound is strictly
above the upper bound in \eqref{eq:truncbound}. The nonzero bounded-support
series $P(T)$ cannot be canceled by the tail. Hence $P(Y)\ne0$.
\end{proof}

\begin{remark}[What the proof does not assume]
Neither $P(T)$ nor the earlier parts $A_j$ need have finite support.
No supremum of their supports in $G$ is required; any displayed upper bound
suffices. There is no rearrangement of a non-summable family and no appeal
to a numerical limit. All polynomial expansions are finite, and each
coefficientwise product is a Hahn convolution.
\end{remark}

\begin{corollary}[A one-series version]\label{cor:onegap}
Let $(a_n)$ be positive and strictly increasing, and suppose that for every
$d\ge1$ and $w\in G$, eventually $a_{n+1}>d a_n+w$. If every $c_n\in K$
is nonzero, then $Y=\sum_n c_nt^{a_n}$ is transcendental over $\FF KG$.
\end{corollary}
\begin{proof}
For a putative relation clear denominators to obtain a nonzero polynomial
$P\in\BB KG[X]$. The case of constant $P$ is impossible. The same tail
estimate as \eqref{eq:tailbound} implies $P(T_n)=0$ for every sufficiently
large $n$, since otherwise its bounded support could not be canceled.
Here $T_n$ is the finite truncation through $n$, supported at most at $a_n$.
These truncations are pairwise distinct, whereas a nonzero polynomial over
a field has only finitely many roots. This contradiction proves the claim.
\end{proof}

The multivariable argument needs more than this root-count observation.
One must ensure that the top homogeneous polynomial does not vanish on all
chosen coefficient vectors. The next section supplies exactly that
certificate, without any algebraic independence assumption on the
coefficients themselves.

\section{Independent coefficient codes and the sharp theorem}\label{sec:independence}

\subsection{A finite-pattern coding lemma}

The following elementary construction is an independent-family coding
argument. Its role is to realize any prescribed finite tuple of positive
integers at cofinally many positions. It is not being claimed as a new
set-theoretic principle.

\begin{lemma}[Cofinal finite-pattern realization]\label{lem:coding}
For any infinite cardinal $\kappa$, there are functions
\[
 c_A:\kappa\longrightarrow\N_{>0}\qquad(A\subseteq\kappa)
\]
such that for every finite list of distinct subsets $A_1,\ldots,A_m$ and
every $(q_1,\ldots,q_m)\in\N_{>0}^m$, the set
\[
 \{\alpha<\kappa:c_{A_j}(\alpha)=q_j\text{ for all }j\}
\]
has cardinality $\kappa$, and in particular is cofinal in $\kappa$.
\end{lemma}
\begin{proof}
Let $D_\kappa$ consist of triples $(s,h,\rho)$ where $s$ is a finite subset
of $\kappa$, $h:\mathcal P(s)\to\N_{>0}$ is a table, and $\rho<\kappa$.
There are $\kappa$ finite subsets of $\kappa$, countably many tables for
each finite $s$, and $\kappa$ choices for $\rho$. Hence
$|D_\kappa|=\kappa$. Fix a bijection
\[
 \alpha\longmapsto(s_\alpha,h_\alpha,\rho_\alpha)
 \quad\text{from }\kappa\text{ to }D_\kappa
\]
and set
\begin{equation}\label{eq:code}
 c_A(\alpha)=h_\alpha(A\cap s_\alpha).
\end{equation}
For each pair $A_j\ne A_k$, choose a point in their symmetric difference.
The finite set $s$ of these points makes $A_j\cap s$ pairwise distinct.
A table $h$ can therefore take the prescribed value $q_j$ at each of them;
set its remaining values to $1$. Every triple $(s,h,\rho)$ realizes the
target tuple, independently of $\rho$. There are $\kappa$ such triples,
which proves the cardinality assertion. A bounded subset of the initial
ordinal $\kappa$ has cardinality smaller than $\kappa$, so the realizing
indices are cofinal.
\end{proof}

\begin{lemma}[Integer grids detect polynomials]\label{lem:grid}
Let $E$ have characteristic zero. A nonzero polynomial
$Q\in E[X_1,\ldots,X_m]$ is nonzero at some point of $\N_{>0}^m$.
More precisely, if each variable has degree at most $d$, any product of
$m$ sets of $d+1$ distinct elements of $E$ contains a nonzero value.
\end{lemma}
\begin{proof}
For one variable, a nonzero polynomial of degree at most $d$ has at most
$d$ roots. For the induction step, regard $Q$ as a polynomial in its last
variable. If it vanishes at every point of the product grid, then for each
fixed tuple of the first $m-1$ variables all its coefficients vanish, by
the one-variable statement. The induction hypothesis applied to those
coefficient polynomials makes each identically zero, a contradiction.
The positive integers are distinct in a characteristic-zero field.
\end{proof}

\subsection{Independence over all bounded-support fractions}

\begin{theorem}[The sharp support theorem]\label{thm:independent}
Let $K$ have characteristic zero and let $0\ne G$ be an ordered abelian
group. Put $\kappa=\cof(G)$. Choose a separated exponent system as in
Lemma~\ref{lem:skeleton} and a coefficient code as in Lemma~\ref{lem:coding}.
Then
\begin{equation}\label{eq:mainfamily}
 Y_A=\sum_{\alpha<\kappa}c_A(\alpha)t^{a_\alpha}
 \qquad(A\subseteq\kappa)
\end{equation}
are algebraically independent over $\FF KG$. Consequently
\begin{equation}\label{eq:lowertrdeg}
 \trdeg_{\FF KG}\HH KG\ \ge\ 2^\kappa.
\end{equation}
Every member has positive integer coefficients and exactly the common
support $\{a_\alpha:\alpha<\kappa\}$, of order type $\kappa$.
\end{theorem}
\begin{proof}
Fix distinct $A_1,\ldots,A_m$ and a nonzero polynomial over $\FF KG$.
Clear a common nonzero denominator to obtain a nonzero polynomial
$P\in\BB KG[X_1,\ldots,X_m]$. If its total degree is zero, its evaluation
is already nonzero. Otherwise let its degree be $d\ge1$ and choose common
coefficient support bounds $[\ell,u]$.

Its top homogeneous part $P_d$ is a nonzero polynomial over the field
$\HH KG$. By Lemma~\ref{lem:grid}, there is a tuple of positive integers
$q=(q_1,\ldots,q_m)$ such that $P_d(q)\ne0$. Lemma~\ref{lem:coding}
realizes this exact tuple at cofinally many indices $\alpha$. The gap
inequalities \eqref{eq:twogaps} hold at every sufficiently large index.
Choose a realizing index after that threshold. Lemma~\ref{lem:detection}
then gives
\[
 P(Y_{A_1},\ldots,Y_{A_m})\ne0.
\]
This handles every finite subfamily and every nonzero polynomial, which is
algebraic independence. There are $2^\kappa$ subsets $A$, proving the lower
bound. Finally the coefficients are positive integers, hence nonzero in
$K$, so none of the designated exponents disappears.
\end{proof}

\begin{corollary}[Optimal cardinality and order type]\label{cor:optimal}
Among all $z\in\HH KG$ transcendental over $\FF KG$, the least support
cardinality is $\cof(G)$ and the least support order type is the initial
ordinal $\cof(G)$.
\end{corollary}
\begin{proof}
A support of cardinality less than $\kappa=\cof(G)$ is bounded above by
definition of cofinality, so its series already belongs to $\BB KG$.
Theorem~\ref{thm:independent} supplies a transcendental series with support
of cardinality and order type exactly $\kappa$. An ordinal smaller than
the initial ordinal $\kappa$ has smaller cardinality, proving the order-type
assertion as well.
\end{proof}

Thus countable supports suffice exactly when the value group has countable
cofinality. When the cofinality is uncountable, every countable-support
series is already a coefficient in the base, regardless of how sparse its
support looks inside a smaller subgroup.

\begin{corollary}[Arbitrary bounded perturbations]\label{cor:perturb}
For any choices $b_A\in\FF KG$ and $u_A\in\FF KG^\times$, the family
$(u_AY_A+b_A)_{A\subseteq\kappa}$ is algebraically independent over
$\FF KG$. In particular arbitrary bounded-support changes, or deletion
of an arbitrary bounded initial part of each series, preserve independence.
\end{corollary}
\begin{proof}
On each finite subfamily the substitutions
$X_A\mapsto u_AX_A+b_A$ define an invertible polynomial change of variables
over the base field. They therefore send nonzero polynomials to nonzero
polynomials. A deleted bounded initial part belongs to $\BB KG$.
\end{proof}

\begin{corollary}[Independent families in every valuation ball]\label{cor:balls}
Every nonempty ball
\[
 U(f,\delta)=\{z\in\HH KG:v(z-f)>\delta\}
\]
contains $2^\kappa$ elements algebraically independent over $\FF KG$.
It also contains an element of $\BB KG$.
\end{corollary}
\begin{proof}
Put $b=f_{\le\delta}\in\BB KG$. Choose $\beta<\kappa$ with
$a_\beta>\delta$, and define
\[
 Z_A=b+\sum_{\alpha\ge\beta}c_A(\alpha)t^{a_\alpha}.
\]
Then $v(Z_A-b)>\delta$ and $v(b-f)>\delta$, with the usual convention for
zero, so every $Z_A$ lies in the ball. Each $Z_A$ differs from $Y_A$ by a
bounded-support series: the omitted initial part is bounded above by
$a_\beta$, and $b$ is bounded. Corollary~\ref{cor:perturb} proves their
independence. The element $b$ itself belongs to the same ball.
\end{proof}

\begin{remark}[A precise limitation of bounded truncation data]
No single bounded truncation can determine algebraicity over $\FF KG$.
The same truncation is shared by a base-field element and by every member
of a family as in Corollary~\ref{cor:balls}. This is an information-theoretic
statement about those data, not a claim that every representation of Hahn
series has an undecidable algebraicity problem.
\end{remark}

\subsection{Exact transcendence degrees}

\begin{corollary}[Cardinality-matched groups]\label{cor:exactcard}
Suppose $\kappa=\cof(G)=|G|$ and $\abs K\le2^\kappa$, with $K$ of
characteristic zero. Then
\[
 \trdeg_{\FF KG}\HH KG=2^\kappa.
\]
\end{corollary}
\begin{proof}
The lower bound is Theorem~\ref{thm:independent}. A Hahn series is a function
from $G$ to $K$, so
\[
 |\HH KG|\le |K|^{|G|}\le(2^\kappa)^\kappa=2^\kappa.
\]
The cardinality of a field bounds its transcendence degree, proving the
upper bound independently of the construction.
\end{proof}

\begin{corollary}[Every nonzero real-exponent subgroup]\label{cor:realgroup}
Let $0\ne G\subseteq\R$ and let $K$ have characteristic zero and
$|K|\le\cc=2^{\aleph_0}$. Then
\begin{equation}\label{eq:realtrdeg}
 \trdeg_{\FF KG}\HH KG=\cc.
\end{equation}
In particular this holds for $K=\Q,\R,\C$ and $G=\Z,\Q,\R$.
\end{corollary}
\begin{proof}
Here $\cof(G)=\aleph_0$, so the lower bound is $\cc$. Every well-ordered
subset $S\subseteq\R$ is countable. Indeed, for each nonmaximal $s\in S$
there is a successor $s^+$ in $S$; choose a rational in $(s,s^+)$. These
intervals are disjoint, producing an injection of $S$ minus a possible
maximum into $\Q$. There are at most $\cc^{\aleph_0}=\cc$ countable
supports in $G$, and at most $|K|^{\aleph_0}\le\cc$ coefficient assignments
to any one support. Thus $|\HH KG|\le\cc$.
\end{proof}

\begin{remark}[Why the full strength is not a counting argument]
For $G=\R$ and $K=\R$ or $\C$, the base $\FF KG$ already has cardinality
$\cc$, since it contains $K$ and lies in a field of cardinality $\cc$.
Equality of these cardinalities gives no lower bound at all on relative
transcendence degree. The separation and coding proof is what supplies it.
Nor is an independent family of maximum possible cardinality automatically
a transcendence basis: maximal cardinality need not mean maximality under
inclusion.
\end{remark}

\section{Explicit countable witnesses and examples}\label{sec:examples}

The coefficient code gives a family of optimal cardinality with integer
coefficients, but its indexing uses a set-theoretic bijection. It is useful
to have additional closed-form witnesses.

\subsection{Factorial gaps with exponential coefficients}

\begin{theorem}[Explicit multiplicative-parameter family]\label{thm:explicit}
Let $u$ be an order unit of $G$ and let $L$ contain $\R$. If
$q_1,\ldots,q_m>0$ are multiplicatively independent real numbers, then
\begin{equation}\label{eq:qfamily}
 Z_{q_j}=\sum_{n\ge2}q_j^n t^{n!u}\qquad(1\le j\le m)
\end{equation}
are algebraically independent over $\FF LG$. Consequently the series
$Z_p$, indexed by the ordinary prime numbers $p$, are jointly algebraically
independent, with positive integer coefficients.
\end{theorem}
\begin{proof}
Multiplicative independence means that
$\prod_jq_j^{m_j}=1$ with $m_j\in\Z$ forces every $m_j=0$.
Let $P\in\BB LG[X]$ be a nonzero polynomial of positive total degree $d$.
At the vector $c(n)=(q_1^n,\ldots,q_m^n)$ its top part is
\[
 P_d(c(n))=\sum_{|\nu|=d}p_\nu Q_\nu^n,
 \qquad Q_\nu=\prod_jq_j^{\nu_j}.
\]
The $Q_\nu$ are distinct positive reals. Choose an exponent $g$ for which
some coefficient $(p_\nu)_g$ is nonzero, and then an $\R$-linear functional
on the finite-dimensional real span of those coefficients in $L$ which is
nonzero on at least one of them. The resulting scalar expression is
\[
 \sum_\nu r_\nu Q_\nu^n,
 \qquad r_\nu\in\R,
\]
with some $r_\nu\ne0$. Divide by the largest $Q_\nu^n$ with nonzero
coefficient. The quotient tends to that nonzero coefficient, so the scalar
expression, and hence $P_d(c(n))$, is nonzero for all sufficiently large
$n$. Only this ordinary finite-dimensional real limit is used, not a Hahn
limit. The factorial gaps satisfy \eqref{eq:twogaps} for large $n$.
Lemma~\ref{lem:detection} proves that $P$ does not vanish on the series.
Clear denominators for a polynomial over the fraction field. Finally,
unique factorization of ordinary integers makes any finite set of distinct
primes multiplicatively independent.
\end{proof}

For $L=\C$ the real-linear functional in this proof can simply be an
appropriate real or imaginary coefficient functional. An alternative proof
uses Corollary~\ref{cor:onescale} to reduce to ordinary Laurent series over
$\R(t^u)$. The direct support proof also records the exact obstruction to
cancellation over bounded-support Hahn coefficients.

\subsection{A continuum family with only integer coefficients}

For $\kappa=\aleph_0$ one can choose an effective enumeration of all triples
\[
 (s,h,\rho),\qquad s\subseteq\N\text{ finite},\quad
 h:\mathcal P(s)\to\N_{>0},\quad\rho\in\N.
\]
Thus for each binary sequence $A\subseteq\N$, formula \eqref{eq:code}
gives an integer coefficient stream. Relabeling its positions by $n\ge2$,
\begin{equation}\label{eq:integerfactorial}
 W_A=\sum_{n\ge2}c_A(n)t^{n!u}
\end{equation}
is a continuum-sized independent family whenever $u$ is an order unit.
Each individual stream is computable relative to the characteristic
function of $A$: evaluating a coefficient asks only finitely many membership
questions. The family is not asserted to consist of computable streams;
there are only countably many computable descriptions.

More strikingly, for $u=1$ these witnesses lie in $\Q[[t]]$, yet they remain
independent over $\FF{\C}{\R}$. All ordinary complex constants and all
bounded real-exponent Hahn series have already been admitted in that base.
Cofinal descent explains why this does not erase ordinary formal-series
transcendence.

\subsection{A small catalogue of contrasting behavior}

The examples below keep the exponent group explicit.

\begin{center}
\small
\begin{tabularx}{\textwidth}{@{}p{0.25\textwidth}X@{}}
\toprule
Series or setting & Consequence \\
\midrule
$\sum_{n\ge0}t^n$, $G=\R$ & Equals $(1-t)^{-1}$ in $\FF{\C}{\R}$, but does not have bounded support. \\
$\sqrt{1+t}$, $G=\R$ & Algebraic of degree two over $\FF{\C}{\R}$; not a bounded-support fraction. \\
$\sum_{n\ge2}t^{n!}$, $G=\R$ & Transcendental over $\FF{\C}{\R}$ by Corollary~\ref{cor:onegap}. \\
$\sum_{n\ge2}p^n t^{n!}$, $p$ prime & A fully explicit countable independent family over the same base. \\
$\R\hookrightarrow \R\omega+\R\subset\No$ & Every series with real exponents becomes bounded in the larger exponent group. \\
$\cof(G)>\aleph_0$ & Every countable-support Hahn series already belongs to $\BB KG$. \\
\bottomrule
\end{tabularx}
\end{center}

The first three rows show why ``unbounded support'' is not synonymous with
``transcendental.'' The final two show why a transcendence assertion must
name its ambient exponent group and its base field.

\section{Rank structure and the unit group}\label{sec:rank}

This section concerns the coefficient ring itself. The field criterion is
proved from the Hahn field construction. The unit classification additionally
uses the known rank-one upper-support theorem, which is isolated explicitly.

\subsection{An exact field criterion}

For a positive element $a\in G$, define its principal convex subgroup
\[
 C(a)=\{g\in G:\abs g\le na\text{ for some }n\ge1\}.
\]
It is proper precisely when $a$ is not an order unit.

\begin{theorem}[Absence of a highest scale]\label{thm:fieldcriterion}
For any field $K$ and nonzero ordered abelian group $G$, the following are
equivalent:
\begin{enumerate}
\item $G$ has no order unit;
\item $\displaystyle\BB KG=\bigcup_{C\subsetneq G\text{ convex}}\HH KC$;
\item $\BB KG$ is a field.
\end{enumerate}
\end{theorem}
\begin{proof}
Assume there is no order unit. If $f$ has support in $[\ell,u]$, choose
$a>0$ with $a\ge\abs\ell,\abs u$. Then its support lies in the proper
convex subgroup $C(a)$, so $f\in\HH K{C(a)}$. Conversely, a proper convex
subgroup $C$ is bounded above in $G$: choose a positive $g\notin C$.
Every positive $c\in C$ is smaller than $g$, since otherwise convexity
would put $g$ in $C$. Hence every series in $\HH KC$ has bounded support
in $G$.

Convex subgroups of an ordered abelian group are linearly ordered by
inclusion. To see this, if neither of $C,D$ contained the other, choose
positive $c\in C\setminus D$ and $d\in D\setminus C$; whichever is smaller
belongs to the other's convex subgroup, a contradiction. Thus the union
in (2) is a directed union of fields and is a field.

Finally suppose $u>0$ is an order unit. The bounded-support element
$1-t^u$ has inverse in the Hahn field
\[
 (1-t^u)^{-1}=\sum_{n\ge0}t^{nu}.
\]
This is the usual strongly summable geometric identity. Its support is
unbounded because $u$ is an order unit. The inverse is unique in the Hahn
field, so no inverse can belong to $\BB KG$. Thus (3) excludes an order
unit and implies (1).
\end{proof}

\begin{corollary}[Closed coefficient fields at unbounded rank]\label{cor:closedbase}
Suppose $G$ is divisible and has no order unit. If $K$ is real closed, then
$\BB KG$ is real closed. If $K$ is algebraically closed of characteristic
zero, then $\BB KG$ is algebraically closed.
\end{corollary}
\begin{proof}
A convex subgroup of a divisible ordered abelian group is divisible: the
unique $n$th part of an element of the subgroup has no larger absolute
value, so belongs to the subgroup by convexity. The standard Hahn-field
closedness criterion makes each $\HH KC$ in Theorem~\ref{thm:fieldcriterion}
real closed or algebraically closed, respectively; this is the imported
criterion recorded in \cite[Fact~2.1.1]{LM24}.

A polynomial has only finitely many coefficients, so all of them lie in
one member of the directed union. In the algebraically closed case it has
a root in that member. In the real closed case the same argument supplies
roots of odd-degree polynomials and square roots of positive elements,
which characterize real closedness. The order throughout is the inherited
Hahn order.
\end{proof}

These fields can be proper and dense in the full Hahn field. There is no
conflict with algebraic closedness: the remaining extension is
transcendental. Theorem~\ref{thm:independent} measures how large it is.

\subsection{Reduction to the highest Archimedean quotient}

Suppose $G$ has an order unit $u$. Define
\begin{equation}\label{eq:topconvex}
 C=\{g\in G:n\abs g<u\text{ for every positive integer }n\}.
\end{equation}
The subgroup $C$ is the unique maximal proper convex subgroup. Indeed,
every $g\notin C$ is itself an order unit in absolute value: some multiple
of $\abs g$ dominates $u$. A convex subgroup containing such a $g$ must
be all of $G$. It follows that $G/C$ is nonzero Archimedean.

\begin{proposition}[Top-rank presentation]\label{prop:toprank}
If $G$ is divisible and has an order unit, choose an ordered embedding of
$H=G/C$ in $\R$ and a $\Q$-linear section of the quotient map. With
$M=\HH KC$, that choice gives isomorphisms
\begin{equation}\label{eq:iterated}
 \HH KG\cong\HH MH,
 \qquad \BB KG\cong\BB MH.
\end{equation}
The full-field isomorphism is the usual iterated Hahn presentation;
the second is its bounded-support restriction.
\end{proposition}
\begin{proof}
Since $G$ and $C$ are $\Q$-vector spaces, choose a complementary subspace.
Order on the quotient dominates order within $C$, so this identifies $G$
with the lexicographic group $H\oplus C$. Regroup a series according to its
$H$ coordinate. Its projected support is well ordered and each fiber is
well ordered; conversely a well-ordered outer support with well-ordered
fibers gives a well-ordered lexicographic support. Coefficientwise addition
and Hahn convolution agree under this grouping, proving the field
isomorphism. This standard presentation is also recorded in
\cite[Fact~2.4.2]{LM24}.

If the original support is bounded by $g$, its $H$ projection is bounded
by the image of $g$. Conversely, if its outer support has upper bound
$h\in H$, choose $h'>h$. A lift of $h'$ bounds every fiber, regardless of
its internal $C$ support. Thus boundedness is equivalent on the two sides.
\end{proof}

\begin{proposition}[Imported rank-one support theorem]\label{prop:supimport}
If $M$ has characteristic zero and $H\subseteq\R$, then for nonzero
$f,g\in\BB MH$,
\[
 \sup_{\R}\supp(fg)=\sup_{\R}\supp(f)+\sup_{\R}\supp(g).
\]
\end{proposition}
This is the monomial-localized form of
\cite[Proposition~3.5.1]{LM24}. It is not reproved or claimed as new here.
The supremum is taken in $\R$ and need not be an element of $H$.

\begin{corollary}[Units see exactly the highest rank]\label{cor:units}
Suppose $K$ has characteristic zero and $G$ is divisible with an order
unit, and let $C$ be \eqref{eq:topconvex}. A nonzero $f\in\BB KG$ is a
unit if and only if its support lies in one coset of $C$. Equivalently,
its image in \eqref{eq:iterated} is a single outer monomial with an arbitrary
nonzero coefficient from $\HH KC$. In particular,
\[
 1+t^g\in\BB KG^\times\quad\Longleftrightarrow\quad g\in C.
\]
\end{corollary}
\begin{proof}
For a nonzero rank-one bounded series define
\[
 w(f)=\sup_{\R}\supp(f)-\min\supp(f)\ge0.
\]
Proposition~\ref{prop:supimport} and the usual valuation identity give
$w(fg)=w(f)+w(g)$. A unit must have width zero, hence singleton support.
Conversely every nonzero monomial is a unit in the bounded-support ring.
Transfer this statement through \eqref{eq:iterated}. For the last claim,
if $g\ne0$ the support is $\{0,g\}$, which lies in one coset precisely
when $g\in C$. If $g=0$, the element is $2$, a nonzero scalar in
characteristic zero, and the same conclusion holds.
\end{proof}

No arbitrary-rank supremum formula has been assumed. In particular the
rank-one width proof should not be repeated in $G$ itself using a guessed
Dedekind completion: at higher rank such a formula requires additional
structure and need not have the same meaning.

\section{Actual surreal and surcomplex realizations}\label{sec:surreal}

\subsection{The normal-form bridge and its scope}

Let $G$ be a set-sized ordered additive subgroup of $\No$. Conway normal
form supplies the embedding
\begin{equation}\label{eq:conway}
 \iota_G:\HH{\R}G\longrightarrow\No,
 \qquad \sum_g r_gt^g\longmapsto\sum_g r_g\omega^{-g}.
\end{equation}
The reversed sign is essential: increasing Hahn exponents correspond to
decreasing Conway exponents. We use the standard normal-form isomorphism
and its field compatibility, as recorded in \cite[Section~2.3]{LM24} and
developed in \cite{Gonshor}. This is an imported foundational theorem,
not an assertion about convergence of real-indexed sums in the fine
surreal topology.

Applying \eqref{eq:conway} to real and imaginary coefficients gives a field
embedding
\[
 \iota_G^{\C}:\HH{\C}G\longrightarrow\No[i].
\]
Indeed, coefficientwise real and imaginary parts have well-ordered
supports, their union is well ordered, and complex multiplication agrees
with multiplication in the quadratic extension. We henceforth identify
these Hahn fields with their concrete images, but continue to name $G$ in
every bounded-support base.

\subsection{Real-exponent surreal numbers}

Define the concrete subfields
\[
 \mathscr H_\R=\left\{\sum_{g\in S}r_g\omega^{-g}:
       S\subseteq\R\text{ well ordered},\ r_g\in\R\right\},
 \qquad \mathscr H_\C=\mathscr H_\R[i].
\]
Let $\mathscr B_\R$ and $\mathscr B_\C$ require $S$ to be bounded above
\emph{in $\R$}, and put $\mathscr F_\R=\Frac\mathscr B_\R$ and
$\mathscr F_\C=\Frac\mathscr B_\C$.

\begin{theorem}[A continuum of independent actual infinitesimals]\label{thm:actualreal}
There is a family of positive infinitesimal surreal numbers
\[
 \xi_A=\sum_{n\ge2}c_A(n)\omega^{-n!}\qquad(A\subseteq\N),
\]
with positive integer coefficients, algebraically independent even over
$\mathscr F_\C$. Moreover
\[
 \trdeg_{\mathscr F_\R}\mathscr H_\R
   =\trdeg_{\mathscr F_\C}\mathscr H_\C=\cc.
\]
The separate explicitly given family
\[
 \xi_p=\sum_{n\ge2}p^n\omega^{-n!}\qquad(p\text{ an ordinary prime})
\]
is also algebraically independent over $\mathscr F_\C$.
\end{theorem}
\begin{proof}
Use the countable coefficient code and factorial separated system in the
complex-coefficient theorem. Its coefficients happen to be positive
integers, so the witnesses lie in the real subfield. Their leading term
is a positive real multiple of $\omega^{-2}$; the remaining terms have
strictly smaller Conway monomials. Thus they are positive infinitesimals.
The exact degrees follow from Corollary~\ref{cor:realgroup}, and the prime
family from Theorem~\ref{thm:explicit}, all transported by
\eqref{eq:conway}.
\end{proof}

This is not an assertion of transcendence over all surreals or all
surcomplex numbers: the witnesses themselves belong to those proper-class
fields. It is a relative algebraic theorem over explicitly specified
set-sized subfields.

\subsection{Maximal families at every infinite regular cardinal}

Let $\kappa$ be an infinite regular cardinal, regarded as an initial ordinal,
and form the ordered additive subgroup of actual surreal numbers
\begin{equation}\label{eq:Gkappa}
 G_\kappa=\bigoplus_{\alpha<\kappa}\Q\,\omega^\alpha.
\end{equation}
The direct sum means finite rational linear combinations. Comparison is
by the largest exponent with a nonzero coefficient. This group has
cardinality $\kappa$, is divisible, has cofinality $\kappa$, and has no
order unit. For the cofinality assertion, any family of fewer than $\kappa$
finite supports uses a bounded set of indices, by regularity; a larger
basis monomial dominates it. Conversely the $\omega^\alpha$ are cofinal.
The absence of an order unit follows by moving one index past the largest
index occurring in a proposed unit.

Let $\mathscr H_\kappa^{\R}$ and $\mathscr H_\kappa^{\C}$ be the real and
complex Hahn fields on this concrete exponent group, interpreted in
$\No$ and $\No[i]$. Let $\mathscr B_\kappa^{\R}$ and
$\mathscr B_\kappa^{\C}$ be their bounded-support rings.

\begin{theorem}[Arbitrarily large, explicit-support surreal families]\label{thm:kappa}
For every infinite regular cardinal $\kappa$:
\begin{enumerate}
\item $\mathscr B_\kappa^{\R}$ is a real closed field and
$\mathscr B_\kappa^{\C}$ is an algebraically closed field.
\item There are $2^\kappa$ positive infinitesimal surreal numbers
\begin{equation}\label{eq:kappaxi}
 \xi_A=\sum_{\alpha<\kappa}
     c_A(\alpha)\,\omega^{-\omega^{\alpha+1}}
     \qquad(A\subseteq\kappa),
\end{equation}
all with positive integer coefficients and exactly the displayed common
support, algebraically independent over $\mathscr B_\kappa^{\C}$.
\item Both exact relative transcendence degrees are $2^\kappa$:
\[
 \trdeg_{\mathscr B_\kappa^{\R}}\mathscr H_\kappa^{\R}
 =\trdeg_{\mathscr B_\kappa^{\C}}\mathscr H_\kappa^{\C}
 =2^\kappa.
\]
No element with a support of smaller cardinality is transcendental over
its corresponding bounded-support base.
\end{enumerate}
\end{theorem}
\begin{proof}
The first assertion is Corollary~\ref{cor:closedbase}. For the second, the
concrete choices
\[
 a_\alpha=\omega^{\alpha+1},\qquad b_\alpha=\omega^\alpha
 \quad(\alpha<\kappa)
\]
form a separated exponent system in $G_\kappa$. Earlier $a_\beta$ are at
most $b_\alpha$. For every finite $n$, $a_\alpha>n b_\alpha$, and the
$b_\alpha$ are cofinal. Thus the two inequalities in \eqref{eq:gaps}
follow exactly as in the uncountable construction, and also hold when
$\kappa=\aleph_0$. Apply Theorem~\ref{thm:independent} over $\C$ and then
the Conway embedding. Since the base ring is already a field, its fraction
field is itself. The leading term of each number is a positive integer
multiple of $\omega^{-\omega}$, proving positivity and infinitesimality.

Finally $|G_\kappa|=\cof(G_\kappa)=\kappa$, while
$|\R|=|\C|=\cc\le2^\kappa$. Corollary~\ref{cor:exactcard} supplies the
upper bounds as well as the equality. Optimality is
Corollary~\ref{cor:optimal}.
\end{proof}

No continuum hypothesis or generalized continuum hypothesis is used. The
regularity assumption on $\kappa$ in this concrete example is genuine:
for a singular initial ordinal the direct sum in \eqref{eq:Gkappa} has
cofinality $\cof(\kappa)$, not $\kappa$. The general theorem still applies
to that group with its actual cofinality, but it does not give the displayed
$2^\kappa$ equality by the same argument.

\subsection{Enlarging a workspace can erase the base distinction}

Consider
\[
 G=\R\omega+\R\subset\No,\qquad H=\R\subset G.
\]
Every real number is smaller than $\omega$ in this ordered group, so $H$
is not cofinal. All real-exponent Hahn series therefore have support
bounded above in $G$. In particular every witness of
Theorem~\ref{thm:actualreal} belongs to $\BB{\C}G$, even though it was
transcendental over $\mathscr F_\C$ in the smaller workspace. This is the
concrete collapse in \eqref{eq:collapse}, not a failure of field embedding.

At the full class level every set of surreal exponents is bounded above by
a surreal number. Consequently the naive class-wide version of
``bounded support in $\No$'' imposes no restriction on a Conway normal form.
All the nontrivial statements in this article are deliberately
workspace-relative. A set-sized cofinal support in $G_\kappa$ is not a
cofinal subclass of all of $\No$.

The valuation topology used for density is likewise the topology of the
named Hahn workspace. The partial sums in the constructions converge there
because their exponents are cofinal in that workspace. This is not a
claim of convergence of a non-eventually-constant set-indexed net in the
full surreal fine topology. The repository's foundational distinction
between those topologies must be retained \cite{RepoCatalogue}.

\section{Verification, formalization, and remaining questions}\label{sec:verification}

\subsection{What was checked computationally}

The accompanying \texttt{verify.py} uses exact rational arithmetic and sparse
finite series. It checks the finite convolution identities underlying
coset projection, coefficient-functional descent, the auxiliary-polynomial
expansion, support-band separation, and the tail bound. It also checks
finite-pattern coding, integer-grid detection, factorial-gap inequalities,
and explicit multiplicative-parameter examples. Its seed, counts, and
outputs are recorded in \texttt{verification.json}. The delivered run passed
3,163 exact assertions in twelve groups. This is an assertion count, not a
count of independently proved theorems.

These checks are regression tests of the finite mechanisms. They are not
proofs of the arbitrary-support, transfinite, or cardinal assertions.
In particular, a finite enumeration does not verify that a coding fiber
has cardinality $\kappa$, and a finite truncation does not verify algebraic
independence of an infinite family. Those assertions are established by
the proofs above. The script uses only the Python standard library and
can be run with
\begin{center}
\texttt{python verify.py --output verification.json}.
\end{center}
The delivered verification record is separate from the PDF build record.

\subsection{A feasible formalization decomposition}

No Lean proof is supplied or claimed. The results nevertheless break into
clear interfaces, without a need to formalize all surreal numbers first.
The first interface is Hahn subseries extraction and its module identity
\eqref{eq:projection}, followed by coefficient-linear maps
\eqref{eq:lambda}. The second is a bounded-support subring with finite
support bounds and an explicit fraction-field embedding. The third is a
finite multivariable polynomial expansion lemma giving
\eqref{eq:bands}--\eqref{eq:tailbound}. These are the main algebraic proof
obligations.

Only after those interfaces are complete is it necessary to add the
cofinality construction, finite-pattern coding, and cardinal estimates.
The actual-surreal consequences are then transported through the
Conway-normal-form embedding. Such a transport must be linked to the
repository's exact proved normal-form bridge, not merely to an abstract
Hahn field with the same informal description. The optional GCD descent
in Appendix~\ref{sec:gcd} belongs to a separate dependency branch.

\subsection{Questions left by these theorems}

The results do not determine
$\trdeg_{\FF KG}\HH KG$ for every large ordered group: the lower bound
$2^{\cof(G)}$ can be smaller than the cardinality of the Hahn field, and no
general matching upper bound is asserted. For coefficient fields of finite
characteristic, the descent and field criterion still hold, and the
independent-family argument works with a countably infinite alphabet in an
infinite coefficient field. The positive-integer formulation is specific
to characteristic zero. Finite coefficient fields require a different
coding argument and are not covered by the independence theorem here.

A second question is whether stronger naturality can be imposed on the
$2^\kappa$-sized independent family. The existence proof chooses an
enumeration of finite tables. It is uniform once that enumeration is
fixed, but it is not invariant under every automorphism of the exponent
group. The explicit prime family removes that choice at the cost of being
only countable.

Finally, the GCD question below concerns divisibility inside the bounded
ring, while the main theorem concerns algebraic independence over its
fraction field. Neither assertion substitutes for the other. In
particular, proving a ring to be GCD would not make the huge transcendental
extension disappear.

\appendix
\section{Normalized GCD descent and a published question}\label{sec:gcd}

This appendix is logically independent of the transcendence theorems. It
provides a full proof of a conditional reduction rather than importing a
recent claimed result as an established premise.

\subsection{The input and the conclusion}

For a characteristic-zero field $K$, let $\overline K$ be an algebraic
closure. Consider the hypothesis
\begin{equation}\tag{GCD$_{\R,\overline K}$}\label{hyp:gcd}
 \BB{\overline K}{\R}\text{ is a GCD domain}.
\end{equation}
A GCD domain is an integral domain in which every pair has a greatest common
divisor, understood up to units. The existence of a GCD for a finite
nonempty family follows by induction. By
Proposition~\ref{prop:supimport}, the units of the rank-one bounded-support
ring are exactly the nonzero monomials.

For a nonzero bounded series define
\begin{equation}\label{eq:normalize}
 N(f)=\lc(f)^{-1}t^{-v(f)}f.
\end{equation}
It has least exponent zero and coefficient $1$ there. Two nonzero bounded
series are associates precisely when their normalizations agree.

\begin{theorem}[Normalized GCD descent]\label{thm:gcddescent}
Assume \eqref{hyp:gcd}. For every nonzero ordered subgroup $H\subseteq\R$,
$\BB KH$ is a GCD domain. More precisely, the normalized GCD in
$\BB{\overline K}{\R}$ of any finite family from $\BB KH$, not all zero,
belongs to $\BB KH$ and is a GCD there. No divisibility or completeness
assumption on $H$ is required.
\end{theorem}

\subsection{Characters separate exponent cosets}

\begin{lemma}\label{lem:characters}
Let $A$ be an abelian group and $L$ an algebraically closed field of
characteristic zero. For every $0\ne a\in A$, there is a character
$\chi:A\to L^\times$ with $\chi(a)\ne1$.
\end{lemma}
\begin{proof}
On the cyclic group generated by $a$, choose a primitive $m$th root of
unity as the image of $a$ if it has finite order $m>1$, and choose $2$ if
it has infinite order. In either case the resulting character separates
$a$ from zero.

We explain why this character extends. The abelian group $L^\times$ is
divisible, because every nonzero scalar has an $n$th root. Suppose a
character has been defined on a subgroup $D\subseteq A$, and choose
$x\notin D$. If no nonzero integer multiple of $x$ lies in $D$, give $x$
any nonzero image. Otherwise the integers $n$ with $nx\in D$ form $m\Z$
for a least positive $m$. Choose $y\in L^\times$ with
$y^m=\chi(mx)$, and set
\[
 \widetilde\chi(d+nx)=\chi(d)y^n.
\]
The choice of $y$ makes this well defined. Thus a character on a proper
subgroup can always be extended to a larger subgroup. The maximal
extension supplied by Zorn's lemma is defined on all of $A$.
\end{proof}

\begin{proof}[Proof of Theorem~\ref{thm:gcddescent}]
Let $d$ be the normalized ambient GCD of $f_1,\ldots,f_m\in\BB KH$.
For every character $\chi:\R/H\to\overline K^\times$, define
\[
 T_\chi\left(\sum_r a_rt^r\right)
       =\sum_r a_r\chi(r+H)t^r.
\]
This is an automorphism of the Hahn field and of its bounded-support
subring. It preserves support, multiplication, and normalization, since
$\chi(0+H)=1$. It fixes every $f_j$. Consequently $T_\chi(d)$ is another
normalized GCD of the same family. Uniqueness of normalization implies
$T_\chi(d)=d$.

If an exponent $r\notin H$ occurred in $d$, Lemma~\ref{lem:characters}
would supply a character with $\chi(r+H)\ne1$. Its nonzero coefficient
would then change, contradicting the equality. Hence
$\supp(d)\subseteq H$.

Next every automorphism in $\Aut(\overline K/K)$ acts coefficientwise and
again preserves normalization and fixes the input family. Thus it fixes
every coefficient of $d$. In characteristic zero the fixed field of all
these automorphisms is $K$. Indeed, an algebraic element outside $K$ has
a separable minimal polynomial with another conjugate; an embedding taking
it to that conjugate extends to an automorphism of the algebraic closure.
Therefore every coefficient of $d$ lies in $K$. Its support is bounded in
$\R$ and lies in the nonzero subgroup $H$, which is cofinal in $\R$;
it is bounded in $H$. Thus $d\in\BB KH$.

Each quotient $f_j/d$ belongs to $\HH KH$ because this is a field. It is
bounded in the ambient ring by divisibility of the ambient GCD, so
cofinality makes it bounded in $H$. Hence $d$ divides every $f_j$ in the
base ring. If $e\in\BB KH$ is any common divisor, ambient universality
gives $e\mid d$ there. The quotient $d/e$ is again in $\HH KH$ and ambient
bounded, hence in $\BB KH$. This proves base-ring universality. The
all-zero family can separately be assigned GCD zero.
\end{proof}

\subsection{Relation to the pre-Schreier question}

L'Innocente and Mantova ask, immediately after their Example~9.1.3, whether
$\BB KH$ is pre-Schreier for a divisible Archimedean group $H$
\cite{LM24}. Since a GCD domain is pre-Schreier, Theorem~\ref{thm:gcddescent}
gives a positive answer \emph{conditional on} \eqref{hyp:gcd}.
The latter hypothesis also follows if the one-sided ring
$\overline K((t^{\R^{\le0}}))$ is a GCD domain, since
Proposition~\ref{prop:localdensity} identifies its monomial localization
with $\BB{\overline K}{\R}$, and localization preserves the GCD property.
These two general commutative-algebra implications are standard; the new
argument in this appendix is the normalized character descent.

There is a further conditional consequence at arbitrary divisible rank.
Suppose the real-exponent GCD input holds over the algebraic closure of
\emph{every} characteristic-zero coefficient field. If $G$ has no order
unit, its bounded-support ring is already a field. If it has an order
unit, Proposition~\ref{prop:toprank} writes its bounded ring as
$\BB MH$ with $H\subseteq\R$ and $M=\HH KC$. Applying the descent theorem
to this coefficient field proves that $\BB KG$ is GCD. Thus a universal
rank-one input would settle bounded-support GCD arithmetic at every
divisible rank.

\begin{warning}[Recent external claim, not a premise of the main results]
A public 2026 project by Dan Abramov, \texttt{gaearon/conway-refinement},
advertises a Lean proof that the one-sided real-exponent Hahn ring is GCD
\cite{Abramov}. Its standalone statement
\texttt{HahnSeriesGCD.lean} has a characteristic-zero hypothesis and the
relevant nonpositive-support definition. The project itself presents its
work as provisional. This audit read the statement and README but did not
build or verify the full proof dependency chain or independently validate
its semantics. Accordingly this article retains \eqref{hyp:gcd} as an
explicit hypothesis. It does not claim an independent resolution of
Conway's refinement conjecture or an unconditional resolution of the
published bounded-support pre-Schreier question.
\end{warning}

\section{Detailed audit and dependency ledger}\label{sec:audit}

\subsection{Novelty candidates versus imported facts}

\begin{longtable}{@{}>{\raggedright\arraybackslash}p{0.28\textwidth}>{\raggedright\arraybackslash}p{0.34\textwidth}>{\raggedright\arraybackslash}p{0.31\textwidth}@{}}
\toprule
Statement & Status in this draft & Dependency boundary \\
\midrule
\endfirsthead
\toprule
Statement & Status in this draft & Dependency boundary \\
\midrule
\endhead
Cofinal coefficient--exponent descent
& Proposed contribution; full proof in Theorem~\ref{thm:descent}.
& Coset extraction, a coefficient-linear functional, finite convolution. \\
Exact cofinality converse
& Proposed companion; Sections~\ref{sec:descent} and \ref{sec:independence}.
& Noncofinal supports become bounded; properness follows from independence. \\
$2^{\cof(G)}$ independent family on optimal support
& Principal proposed contribution; Theorem~\ref{thm:independent}.
& Full support-gap and coding proofs; no GCD input. \\
Earlier nonrationality/cofinality obstruction
& Prior result, not claimed new.
& L'Innocente--Mantova, Proposition~2.4.5. \\
Finite-pattern coding and integer-grid detection
& Elementary background mechanisms, proved here.
& Choice and finite polynomial root bounds. \\
Real-exponent degree $\cc$ and arbitrary-regular-cardinal surreal examples
& Consequences of the principal theorem, with explicit constructions.
& Independent cardinal upper bounds; standard normal form for actual surreals. \\
Field criterion and top-rank unit formulation
& Structural companion results; modest novelty claim.
& Hahn fields; top-rank grouping is standard; rank-one upper-support theorem imported. \\
Hahn real/algebraic closedness and Conway normal form
& Imported foundations, not new proofs.
& Standard results explicitly identified in the text. \\
Normalized GCD descent
& Proposed conditional reduction; full implication proof.
& Existence of the ambient GCD remains a hypothesis. \\
Lean certification of this article
& Not claimed.
& Finite exact tests are not a proof-assistant check. \\
\bottomrule
\end{longtable}

\subsection{Repository comparison actually performed}

The repository's root overview, documentation catalogue, current
\texttt{docs/new} tree, and the tail-span report's detailed README were
read through the GitHub connector. The initial inspected revision was
\auditcommit. The \texttt{docs/new} tree at that revision contained its
README rather than an additional unreviewed article collection. Targeted
repository searches for \texttt{lacunary} and \texttt{bounded support}
returned no indexed matches; that negative search result alone is not
proof of absence. The comparison relies principally on the catalogue and
the closest report's explicitly described base fields and conclusions.

The audit was not a line-by-line comparison against all 36 full manuscripts,
all retired sources in Git history, or all Lean declarations. It therefore
supports the formulation ``not identified in the reviewed material,'' not
a guaranteed absence from every historical file. The delivered
\texttt{research\_audit.md} records this limitation alongside the literature
queries and the exact mathematical strengthening over the closest published
cofinality result.

\subsection{Proof audit: the points most vulnerable to mistakes}

The proofs depend on several distinctions that deserve an explicit final
check. The base is the fraction field of bounded-support series, so every
putative relation must first have its denominators cleared. At limit
indices the earlier partial series are not finite; the proof uses bounded
support, not finite support. A top homogeneous polynomial need not be
nonzero on an arbitrary coefficient sequence, which is why the coding
lemma realizes a complete positive-integer grid cofinally. The coding
family has size $2^\kappa$ but uses only $\kappa$ positions; finite patterns,
not arbitrary infinite patterns, are what it realizes.

Coset projection is not a ring homomorphism; it is linear over the Hahn
field on the subgroup. Likewise the coefficient functional is not a field
homomorphism; it is linear over the smaller coefficient field. Cofinality
is exactly what makes their output coefficients bounded in the smaller
group. For units, upper-support multiplicativity is used only after
passing to an Archimedean quotient. For concrete surreal examples, the
monomial map is Conway's $x\mapsto\omega^x$, and the Hahn valuation sign
is reversed in \eqref{eq:conway}. Finally, a maximal-cardinality independent
family is not claimed to be a transcendence basis.

\subsection{What the deliverable establishes}

The unconditional core is a collection of relative algebraic theorems
with detailed proofs and explicit witnesses, including the sharp support
threshold, coefficient--exponent descent, and arbitrarily large actual
surreal examples. Their mathematical validity can be assessed directly
from the given arguments and the named standard foundations. Their
publication novelty requires further expert literature review. The
optional GCD reduction has an additional stated hypothesis and should not
be promoted into an unconditional open-problem solution by silently
accepting an unverified external claim.

\begin{thebibliography}{9}
\raggedright
\bibitem{LM24}
Sonia L'Innocente and Vincenzo Mantova,
\emph{A factorisation theory for generalised power series and omnific integers},
Advances in Mathematics \textbf{442} (2024), article 109513.
\href{https://doi.org/10.1016/j.aim.2024.109513}{doi:10.1016/j.aim.2024.109513}.
Version examined: \href{https://arxiv.org/abs/1710.07304v5}{arXiv:1710.07304v5},
22 January 2024. Relevant locations: Fact 2.1.1, Section 2.3,
Fact 2.4.2, Proposition 2.4.5, Proposition 3.5.1, and the question after
Example 9.1.3.

\bibitem{Gonshor}
Harry Gonshor, \emph{An Introduction to the Theory of Surreal Numbers},
London Mathematical Society Lecture Note Series 110, Cambridge University
Press, 1986.
\href{https://doi.org/10.1017/CBO9780511629143}{doi:10.1017/CBO9780511629143}.
Cited for standard normal-form foundations; bibliographic data also
recorded in \cite{LM24}.

\bibitem{StacksLD}
The Stacks Project Authors,
\emph{The Stacks Project}, Section 9.27, ``Linearly disjoint extensions,''
Tag \href{https://stacks.math.columbia.edu/tag/09IC}{09IC}.
Accessed 22 September 2026.

\bibitem{RepoCatalogue}
Vladimir Reshetnikov, \emph{Surreal}, repository documentation catalogue,
\texttt{docs/README.md}, revision
\texttt{048b72cf7cbfc8ab246e4f73788c10460cb3f6e0}.
\href{https://github.com/VladimirReshetnikov/Surreal/blob/048b72cf7cbfc8ab246e4f73788c10460cb3f6e0/docs/README.md}{Versioned repository source}.
Accessed 22 September 2026.

\bibitem{RepoTail}
\emph{Tail-spans and differential transcendence in surreal and surcomplex
Hahn fields}, AI-assisted research report in \repo,
\texttt{docs/surreal/tail-spans-and-differential-transcendence/}.
Detailed README examined at the same revision as \cite{RepoCatalogue}.
\href{https://github.com/VladimirReshetnikov/Surreal/blob/048b72cf7cbfc8ab246e4f73788c10460cb3f6e0/docs/surreal/tail-spans-and-differential-transcendence/README.md}{Versioned report overview}.
Not treated as a refereed source.

\bibitem{Abramov}
Dan Abramov, \emph{A Proof of Conway's Refinement Conjecture},
public provisional project \texttt{gaearon/conway-refinement}, 2026.
\href{https://github.com/gaearon/conway-refinement}{Project README} and
\href{https://github.com/gaearon/conway-refinement/blob/main/ConwayRefinement/Standalone/Mathlib/HahnSeriesGCD.lean}{standalone Hahn-series GCD statement}.
Accessed 22 September 2026. The proof build was not independently rerun in
this audit; the stated GCD input is retained as a hypothesis.
\end{thebibliography}
\end{document}
